\documentclass{subfiles}
\begin{document}
    So far we have considered the case in which the set of keys \(K\) was dynamic.
        We now turn our attention to the case in which \(K\) is static,
        and we show that in this case a constant time lookup worst case is achiveable.

    The scheme we are about to discuss is often refferred to as \emph{perfect hashing}.
        The idea is the following: we use a two level schema in which,
        by means of a hash function \(h\) randomly choosed from a family of universal hash functions,
        each of the \(n\) keys is mapped to one of the \(m\) slots of the table.
        Each slot \(j\) stores a secondary hash table of size \(m_{j} = n_{j}^{2}\) 
        to which is associate a hash function \(h_{j}\), 
        where \(n_{j}\) is the number of keys mapping to \(j\).
    
    Assume that the first level hash function comes from \(\H_{pm}\), 
        while each of the secondary level hash functions come from \(\H_{pm_{j}}\).
        We show that, choosing \(h_{j}\) carefully we can garantee that no collision will occur.
        Additionally, we show that the overall space required by both the primary table 
        and the secondary one is amounts to \(\ofO{n}\).
        The following theorems and corollary will do the work.

    \begin{theorem*}
        Let \(h\) be a universal hash function, and let \(m = n^{2}\) be the size of a hash table 
            used to store a static set of keys \(K\) of size \(n\). Then,
            the probability of a collision occuring is les than \(\sfrac{1}{2}\).
    \end{theorem*}
    \begin{proof}
        Since \(h\) is a universal hash function, any two keys collide with probability \(\sfrac{1}{m}\).
            The number of possible pair that may collide is \(\binom{n}{2}\).
            If \(X\) is a random variable that counts the number of collisions, 
            then 
            \[\begin{aligned}
                \E[X]{} & = \binom{n}{2} \frac{1}{m} \\
                        & = \binom{n}{2} \frac{1}{n^{2}} && \text{ since \(m = n^{2} \)}\\ 
                        & = \frac{n^{2} - n}{2} \frac{1}{n^{2}} \\ 
                        & < \frac{1}{2} && \text{as \(n \to \infty\)}
            \end{aligned}\]
    \end{proof}

    \begin{theorem*}
        Let \(h\) be a hash function randomly chosen from a family of universal hash functions,
            and assume we store \(n\) keys in a table of size \(m = n\).
            Then, it holds
            \[
                \E[\sum_{i = 0}^{m - 1} n_{j}^{2}]{} < 2,
            \]
            \(n_{j}\) being the number of keys hashing to \(j\).
    \end{theorem*}
    \begin{proof}
        By a simple induction we have that the following relation holds for every non-negative integer:
        \[
            k^{2} = k + 2\binom{k}{2}.
        \]

        It follows 
        \[\begin{aligned}
            \E[\sum_{i = 0}^{m - 1} n_{j}^{2}]{} & =
                \E[\sum_{i = 0}^{m - 1} n_{j}]{} +
                2\E[\sum_{i = 0}^{m - 1} \binom{n_{j}}{2}]{} && \text{by the linearity of expectation} \\ 
                & = \E[n]{} + 2\E[\sum_{i = 0}^{m - 1} \binom{n_{j}}{2}]{} \\ 
                & = n + 2\E[\sum_{i = 0}^{m - 1} \binom{n_{j}}{2}]{} && \text{since \(n\) is not a random variable}
        \end{aligned}\]
        Now, recall that \(h\) is a universal hash function, hence 
        \[\begin{aligned}
            \E[\sum_{i = 0}^{m - 1} \binom{n_{j}}{2}]{} & \le \binom{n}{m}\frac{n}{m}  \\ 
                & = \frac{n(n - 1)}{2m} \\
                & = \frac{n - 1}{2}
        \end{aligned}\]
        since \(m = n\). Therefore, 
        \[
            \E[\sum_{i = 0}^{m - 1} n_{j}^{2}]{} \le n + 2 \frac{n - 1}{2} < 2n
        \]
    \end{proof}
    \begin{corollary*}
        Let \(h\) be a hashing function randomly chosen from a family of universal hash functions,
            and assume we store \(n\) keys in an hash table of size \(m = n\),
            and we set the size of each secondary hash table to \(m_{j} = n^{2}_{j}\).
            Then, under a perfect hashing schema, the amount of space required is no more than \(2n\).
    \end{corollary*}
    \begin{proof}
        Follows immediately from the previous theorem.
    \end{proof}

    \nocite{CormenLRS2009}
\end{document}
